% 4-5-1-dm-ours.tex
%  Budget: ~4pp.  HANDOFF 6.3, 6.4, 19.1 (dm_tests, dm_regime_split,
%  dm_bootstrap_sensitivity, oracle_bound), and 16A/A1.
%  Source map: no external literature claims; all values from frozen files.
%  MANDATE (16A/A1):
%   - report raw AND Holm p for the 21-cell exploratory family
%   - LSTM vs LightGBM is NOT significant (raw 0.0405, Holm 0.1215)
%   - regime-aware vs static carries all three numbers together:
%     pooled 0.0226, stressed 0.0063, bootstrap 0.0129-0.0571
%  Rules: decimals in \lr{}; integers Persian; every table cell in \lr{}.

\subsection{مقایسه مدل‌های این پژوهش}
\label{sec:4-5-1}

\subsubsection*{ماتریس کامل مقایسه‌های دوبه‌دو}

هفت مدل \lr{—} پنج مدل منفرد و دو ترکیب \lr{—} ۲۱ مقایسه‌ی دوبه‌دو تولید می‌کنند.
\cref{tab:dm-matrix} همه‌ی آن‌ها را می‌آورد. هیچ سلولی حذف نشده است؛ گزینشِ سلول‌ها
پس از دیدن نتیجه، همان کاری است که تصحیح چندگانگی برای مهارش وجود دارد.

\begin{table}[htbp]
\centering
\caption{\lr{p}-مقدارهای خامِ آزمون دیبولد-ماریانو با تصحیح \lr{HAC}، افق ساعتی،
۷۲۸ مبدأ. هر درایه احتمالِ آن است که مدلِ \emph{سطر} بهتر از مدلِ \emph{ستون}
نباشد؛ مقدار نزدیک به صفر یعنی سطر بهتر است. ستون‌ها به همان ترتیب سطرها
شماره‌گذاری شده‌اند. آنچه می‌آید ستونِ خام است و به‌تنهایی مبنای هیچ ادعای معناداری قرار
نمی‌گیرد \lr{—} \cref{tab:dm-holm} را ببینید.}
\label{tab:dm-matrix}
\small
\setlength{\tabcolsep}{4pt}
\begin{tabular}{lrrrrrrr}
\toprule
مدل & ۱ & ۲ & ۳ & ۴ & ۵ & ۶ & ۷ \\
\midrule
۱ ساده & \lr{—} & \lr{1.000} & \lr{1.000} & \lr{1.000} & \lr{1.000} & \lr{1.000} & \lr{1.000} \\
۲ \lr{SARIMAX} & \lr{0.000} & \lr{—} & \lr{1.000} & \lr{0.9999} & \lr{1.000} & \lr{1.000} & \lr{1.000} \\
۳ \lr{LEAR-LASSO} & \lr{0.000} & \lr{0.000} & \lr{—} & \lr{0.2665} & \lr{0.5964} & \lr{1.000} & \lr{1.000} \\
۴ \lr{LightGBM} & \lr{0.000} & \lr{0.0001} & \lr{0.7335} & \lr{—} & \lr{0.9595} & \lr{1.000} & \lr{1.000} \\
۵ \lr{LSTM} & \lr{0.000} & \lr{0.000} & \lr{0.4036} & \lr{0.0405} & \lr{—} & \lr{1.000} & \lr{1.000} \\
۶ ترکیب ایستا & \lr{0.000} & \lr{0.000} & \lr{0.000} & \lr{0.000} & \lr{0.000} & \lr{—} & \lr{0.9774} \\
۷ ترکیب رژیم‌محور & \lr{0.000} & \lr{0.000} & \lr{0.000} & \lr{0.000} & \lr{0.000} & \lr{0.0226} & \lr{—} \\
\bottomrule
\end{tabular}
\end{table}

\subsubsection*{ماتریس پس از تصحیح هولم\lr{–}بونفرونی}

اِعمال روش هولم\lr{–}بونفرونی بر هر ۲۱ آزمون، ۱۷ آزمون را دست‌نخورده باقی می‌گذارد
و ۴ آزمون را از فهرست نتایج معنادار بیرون می‌برد. \cref{tab:dm-holm} خلاصه‌ی
آن است. هیچ مدلی برای این محاسبه بازاجرا نشد؛ تصحیح صرفاً بر همان
\lr{p}-مقدارهای تثبیت‌شده‌ی \cref{tab:dm-matrix} انجام گرفته است.

\begin{table}[htbp]
\centering
\caption{نتیجه‌ی تصحیح هولم\lr{–}بونفرونی بر خانواده‌ی اکتشافیِ ۲۱ آزمونه.
«مقاوم» یعنی نتیجه پس از تصحیح نیز در سطح \lr{0.05} برقرار می‌ماند.}
\label{tab:dm-holm}
\small
\setlength{\tabcolsep}{5pt}
\begin{tabular}{lrrl}
\toprule
مقایسه & \lr{p} خام & \lr{p} هولم & وضعیت \\
\midrule
هر شش مدل در برابر ساده & \lr{0.000} & \lr{0.000} & مقاوم \\
دو ترکیب در برابر هر عضو (۸ آزمون) & \lr{$\le$4.8e-06} & \lr{$\le$3.4e-05} & مقاوم \\
\lr{SARIMAX} در برابر سه مدل دیگر & \lr{$\le$1.0e-04} & \lr{$\le$5.1e-04} & مقاوم \\
\midrule
ترکیب رژیم‌محور در برابر ایستا & \lr{0.0226} & \lr{0.0903} & مقاوم نیست \\
\lr{LSTM} در برابر \lr{LightGBM} & \lr{0.0405} & \lr{0.1215} & مقاوم نیست \\
\lr{LEAR-LASSO} در برابر \lr{LightGBM} & \lr{0.2665} & \lr{0.5329} & از ابتدا تهی \\
\lr{LEAR-LASSO} در برابر \lr{LSTM} & \lr{0.4036} & \lr{0.5329} & از ابتدا تهی \\
\bottomrule
\end{tabular}
\end{table}

هرچه ساختاری است، دست‌نخورده می‌ماند. اینکه هر چهار مدل آموخته بر مدل ساده برتری
دارند و اینکه هر دو ترکیب بر تک‌تک اعضای خود برتری دارند، پس از تصحیح نیز با
فاصله‌ی زیادی برقرار است. آنچه از دست می‌رود، دو ادعای مرزی است.

\subsubsection*{سه مدل که از هم جدا نمی‌شوند}

\lr{LSTM} با \lr{3.8734}، \lr{LEAR-LASSO} با \lr{3.8988} و \lr{LightGBM} با
\lr{3.9678} در فاصله‌ی کمتر از یک‌دهم یورو بر مگاوات‌ساعت از یکدیگر قرار دارند، و
آزمون هیچ‌کدام را از دیگری جدا نمی‌کند.

مقایسه‌ی \lr{LSTM} با \lr{LEAR-LASSO} حتی پیش از تصحیح نیز تهی است:
\lr{p} برابر \lr{0.4036}. عددی که جای تفسیرِ ظریف باقی نمی‌گذارد. بهترین مدل آماری و بهترین
مدل یادگیری عمیقِ این پژوهش، بر این بازار و این پنجره، از یکدیگر تفکیک‌پذیر نیستند.

مقایسه‌ی \lr{LSTM} با \lr{LightGBM} ظریف‌تر است و باید با دقت نوشته شود.
\lr{p}-مقدار خام \lr{0.0405} است، که در نگاه نخست معنادار به نظر می‌رسد. اما این
آزمون در خانواده‌ی اکتشافیِ ۲۱تایی قرار دارد و پس از تصحیح هولم به \lr{0.1215}
می‌رسد. بنابراین: \lr{LSTM} برآورد نقطه‌ای پایین‌تری دارد، و این تفاوت تصحیح برای
چندگانگی را تاب نمی‌آورد.

این تصمیمِ نوشتاری، برخلاف ظاهرش، پاسخ به پرسش پژوهشی نخست را \emph{تقویت} می‌کند.
خوانشِ صادقانه این می‌شود که سه خانواده‌ی مدلی متفاوت \lr{—} خطیِ منظم‌شده، درختِ
تقویت‌شده و شبکه‌ی بازگشتی \lr{—} با بودجه‌ی تنظیم برابر به دقتی می‌رسند که آزمون
آماری از هم تمیزشان نمی‌دهد؛ و تنها ترکیب است که از این سه فاصله می‌گیرد. ادعای
ضعیف‌تر اما مقاومِ «تفکیک‌ناپذیرند» بسیار دفاع‌پذیرتر از ادعای قوی‌ترِ «\lr{LSTM}
بهتر است» است که با یک تصحیح استاندارد فرو می‌ریزد.

\subsubsection*{ترکیب: تنها جدایی مقاوم}

هر دو ترکیب بر تک‌تک اعضای خود برتری دارند، و این هشت مقایسه پس از تصحیح هولم نیز
با \lr{p} حداکثر \lr{3.4e-05} برقرار می‌مانند. مقاوم‌ترین یافته‌ی این بخش پس از
برتریِ بدیهی بر مدل ساده است.

اندازه‌ی این برتری اما متواضع است. ترکیب ایستا \lr{3.5742} و ترکیب رژیم‌محور
\lr{3.5569} در برابر بهترین عضوِ منفرد با \lr{3.8734}؛ یعنی بهبودی حدود هشت درصد.
\cref{sec:3-8} کرانِ اوراکل را \lr{3.5580} گزارش کرد: بهترین وزن‌های ایستایی که با
دانستنِ کاملِ دوره‌ی آزمون قابل انتخاب بودند. ترکیب ایستای واقعی \lr{3.5742} است، پس
فاصله تا آن کران کمتر از دو صدم یورو بر مگاوات‌ساعت است. عملاً هیچ وزن‌دهیِ ایستایی
نمی‌توانست از این بهتر عمل کند، و این خودْ توضیح می‌دهد که چرا گام بعدی مشروط‌کردن
وزن‌ها به رژیم بود و نه جست‌وجوی وزن‌های بهتر.

\subsubsection*{وزن‌دهی رژیم‌محور: آزمون تأییدی}

تنها فرضیه‌ی خانواده‌ی تأییدی همین است، و باید با هر سه عددش با هم بیان شود.

بر مجموع ۷۲۸ روز، ترکیب رژیم‌محور با \lr{3.5569} در برابر \lr{3.5742} برتری دارد و
\lr{p} برابر \lr{0.0226} است. به‌عنوان یک فرضیه‌ی از پیش ثبت‌شده، این در سطح
\lr{0.05} معنادار است. اما دو قید باید بی‌درنگ کنارش بیاید. نخست آنکه اگر همین آزمون
از خانواده‌ی اکتشافی برداشته شده بود، پس از تصحیح هولم به \lr{0.0903} می‌رسید و
معنادار نمی‌ماند؛ دفاع از آن تماماً بر ثبتِ تاریخ‌دارِ پیش از اجرا استوار است، نه بر
اندازه‌ی خودِ عدد. دوم آنکه، چنان‌که \cref{tab:dm-bootstrap} نشان می‌دهد، تحت
محافظه‌کارانه‌ترین فرض وابستگی، همین \lr{p} از \lr{0.05} عبور می‌کند.

\cref{tab:dm-regime} نشان می‌دهد این بهبود از کجا می‌آید، و پاسخ تیزتر از چیزی است
که میانگین کل نشان می‌دهد.

\begin{table}[htbp]
\centering
\caption{تفکیک برتریِ ترکیب رژیم‌محور بر ترکیب ایستا به تفکیک رژیم بازار.
\lr{MAE} بر حسب یورو بر مگاوات‌ساعت.}
\label{tab:dm-regime}
\small
\setlength{\tabcolsep}{6pt}
\begin{tabular}{lrrrr}
\toprule
زیرمجموعه & روز & رژیم‌محور & ایستا & \lr{p} \\
\midrule
کل دوره & ۷۲۸ & \lr{\textbf{3.5569}} & \lr{3.5742} & \lr{0.0226} \\
روزهای پرتنش & ۷۷ & \lr{\textbf{5.5132}} & \lr{5.6830} & \lr{0.0063} \\
روزهای آرام & ۶۵۱ & \lr{3.3255} & \lr{\textbf{3.3248}} & \lr{0.8465} \\
\bottomrule
\end{tabular}
\end{table}

تمام بهبود از ۷۷ روز پرتنش می‌آید. آنجا اختلاف \lr{0.1698} یورو بر مگاوات‌ساعت و
\lr{p} برابر \lr{0.0063} است. بر ۶۵۱ روز آرام، ترکیب رژیم‌محور به‌اندازه‌ی
ده‌هزارم یورو \emph{بدتر} است و \lr{p} برابر \lr{0.8465}؛ یعنی هیچ.

سازوکار دقیقاً برای انجام همین ساخته شده بود. وزن‌های رژیم آرام و
وزن‌های ایستا تقریباً یکسان‌اند \lr{—} \cref{tab:ensemble-weights} این را نشان
می‌دهد \lr{—} چون ۳۲۰ روز از ۳۵۷ روزِ پنجره‌ی برازش آرام بودند و عملاً همان وزن‌ها را
تعیین کردند. تفاوت تنها در سبد پرتنش پدید می‌آید. پس بیان درست ادعا این است:
مشروط‌کردن وزن‌ها به رژیم، در روزهای پرتنش کمک می‌کند و در روزهای آرام بی‌اثر است.
این ادعایی محدودتر از «ترکیب رژیم‌محور بهتر است» می‌باشد و برخلاف آن، هم با اعداد
سازگار است و هم با هدفِ اعلام‌شده‌ی طراحی.

\subsubsection*{آزمون مقاومت: بوت‌استرپ بلوکی}

آزمون دیبولد-ماریانو با تصحیح \lr{HAC} خودهمبستگیِ سری اختلاف زیان را تا مرتبه‌ای
معین جبران می‌کند. برای سنجش حساسیت نتیجه به این فرض، همان مقایسه با بوت‌استرپ بلوکی
در طول‌های بلوک ۳ تا ۱۰ روز نیز اجرا شد.

\begin{table}[htbp]
\centering
\caption{حساسیت نتیجه‌ی «رژیم‌محور در برابر ایستا» به طول بلوک در بوت‌استرپ بلوکی.
ستون آخر \lr{p}-مقدار \lr{HAC} برای همان زیرمجموعه است و در هر سطر ثابت می‌ماند.}
\label{tab:dm-bootstrap}
\small
\setlength{\tabcolsep}{6pt}
\begin{tabular}{lrrr}
\toprule
زیرمجموعه & طول بلوک & \lr{p} بوت‌استرپ & \lr{p} \lr{HAC} \\
\midrule
کل دوره & ۳ & \lr{0.0129} & \lr{0.0226} \\
کل دوره & ۵ & \lr{0.0269} & \lr{0.0226} \\
کل دوره & ۷ & \lr{0.0396} & \lr{0.0226} \\
کل دوره & ۱۰ & \lr{0.0571} & \lr{0.0226} \\
\midrule
پرتنش & ۳ & \lr{0.0081} & \lr{0.0063} \\
پرتنش & ۵ & \lr{0.0193} & \lr{0.0063} \\
پرتنش & ۷ & \lr{0.0311} & \lr{0.0063} \\
پرتنش & ۱۰ & \lr{0.0439} & \lr{0.0063} \\
\midrule
آرام & ۳ & \lr{0.8525} & \lr{0.8465} \\
آرام & ۱۰ & \lr{0.8399} & \lr{0.8465} \\
\bottomrule
\end{tabular}
\end{table}

نتیجه‌ی این آزمون باید بی‌کم‌وکاست گزارش شود. بر کل دوره، \lr{p} از \lr{0.0129} در
بلوک‌های سه‌روزه تا \lr{0.0571} در بلوک‌های ده‌روزه تغییر می‌کند \lr{—} یعنی زیر
محافظه‌کارانه‌ترین فرضِ وابستگی، از آستانه‌ی \lr{0.05} عبور می‌کند. بر زیرمجموعه‌ی
پرتنش، بازه \lr{0.0081} تا \lr{0.0439} است و در سراسر آن معنادار می‌ماند. بر
روزهای آرام، \lr{p} در حدود \lr{0.84} تا \lr{0.85} ثابت است.

جمع‌بندی صادقانه چنین است: نتیجه‌ی تجمیع‌شده مرزی است و نتیجه‌ی زیرمجموعه‌ی پرتنش
پابرجاست. این دو جمله با هم، همان ادعایی را می‌سازند که این پژوهش می‌تواند از آن
دفاع کند، و ننوشتنِ جمله‌ی نخست، جمله‌ی دوم را نیز مشکوک می‌کرد.
